\section{When Ensemble is Irreplaceable}
\label{sec:when-ensemble}

\subsection{Two Use Cases that Require Ensemble Methods}

Despite the computational advantages of \cs{} for plan generation,
there are two use cases where MCMC ensemble methods are irreplaceable:

\paragraph{Use Case 1: Auditing enacted maps.}
When an enacted map has been drawn by a legislature (or by a commission
that may have introduced bias), the question is whether the map is an
outlier relative to algorithmically neutral alternatives.
\cs{} generates one plan; it cannot answer the question ``how unusual
is this plan?''
An ensemble of plans drawn from the neutral distribution provides the
reference distribution needed for the outlier test.

\paragraph{Use Case 2: Characterising the feasible space.}
Researchers and courts may want to understand the range of outcomes
achievable by neutral redistricting.
Questions like ``what is the minimum possible number of majority-Black
districts?'' or ``what is the typical compactness of random redistricting?''
require sampling from the distribution, not finding one optimum.
\cs{} addresses neither question.

\subsection{The Division of Labor}

We recommend the following division of labor between \cs{} and ensemble
methods in the DIA statutory framework:

\begin{enumerate}
\item \textbf{Statutory plan generation}: \cs{} with $T = 600$.
      One plan, deterministic, certifiably optimal, verifiable in 60 seconds.
\item \textbf{Partisan outlier testing}: Certified \recom{} ensemble
      with G.4 diagnostics.
      Used to challenge enacted plans in state court under the
      \textit{Rucho} framework.
\item \textbf{Research and parameter calibration}: Uncertified \recom{}
      for exploratory analysis, followed by certified runs for publication.
\end{enumerate}

\subsection{What Ensemble Methods Cannot Provide}

Ensemble methods cannot provide:
\begin{itemize}
\item A single ``best'' plan with a certifiable quality guarantee.
\item Determinism (without fixing a seed, which defeats the mixing purpose).
\item A fast answer: even the minimum ensemble for NC requires 5{,}000 steps
      and 20 minutes on consumer hardware.
\end{itemize}

These limitations are not defects; they follow from the nature of
distributional inference.
The appropriate response is to use the right tool for the right task,
not to ask ensemble methods to do what \cs{} was designed for.

\subsection{Hybrid Approach for Special Masters}

For special masters appointed by courts to draw remedial plans, a hybrid
approach is appropriate:
\begin{enumerate}
\item Generate the initial plan using \cs{} (optimal, transparent,
      certifiable).
\item Run a certified ensemble audit to verify that the \cs{} plan is
      not an outlier in a partisan direction.
\item If the \cs{} plan is an outlier (rare---expected frequency $< 1\%$
      by B.7), adjust using a constrained bisection with modified objective.
\end{enumerate}

This hybrid approach provides both the plan-generation benefits of \cs{}
and the distributional-verification benefits of ensemble methods.
